% Reused from chapters/09-validation-preprocessing.tex
\begin{labvisual}{Recall the chapter model}
\begin{labgraphic}
\begin{tikzpicture}[x=\linewidth,y=1cm,
  foldcell/.style={draw=UTPBlue!70,fill=UTPPaleBlue,rounded corners=1mm,
    line width=0.6pt,text width=0.13\linewidth,minimum height=0.62cm,
    inner xsep=1mm,inner ysep=1mm,align=center,
    font=\scriptsize\color{UTPNavy}},
  heldout/.style={foldcell,draw=UTPTeal!80,fill=UTPPaleTeal,
    font=\scriptsize\bfseries\color{UTPNavy}},
  heading/.style={font=\scriptsize\bfseries\color{UTPNavy},inner sep=0pt}]
\node[heading,anchor=east] at (0.095,0.65) {Run};
\foreach \c in {0,...,4}{
  \pgfmathtruncatemacro{\foldnumber}{\c+1}
  \node[heading] at (0.20+\c*0.16,0.65) {Fold \foldnumber};
}
\foreach \r in {0,...,4}{
  \pgfmathtruncatemacro{\runnumber}{\r+1}
  \node[heading,anchor=east] at (0.095,-\r*0.78) {\runnumber};
  \foreach \c in {0,...,4}{
    \ifnum\r=\c
      \node[heldout] at (0.20+\c*0.16,-\r*0.78) {Validate};
    \else
      \node[foldcell] at (0.20+\c*0.16,-\r*0.78) {Train};
    \fi
  }
}
\node[font=\scriptsize\color{UTPMuted},align=center,
  text width=0.86\linewidth] at (0.50,-3.85)
  {Each run uses four training folds and one held-out validation fold.};
\end{tikzpicture}
\end{labgraphic}\end{labvisual}
